% !TeX root = ../初等数学笔记.tex
% 由主文档通过 \input 载入；请编译主文档。
\setcounter{chapter}{0} % 保留原书章号。
\chapter{集合、逻辑与证明}
\label{chap:em-language}
数学中的困难往往不只在计算，还在于条件是否清楚、推理是否有效。集合描述研究范围，逻辑说明条件之间的关系，证明则把已知事实连接到待证结论。本章采用朴素集合语言，不讨论集合论的公理化基础。

\section{集合、元素与常用表示}
\begin{definition}
\emph{集合}（set）是由确定对象组成的整体，对象称为元素。$x\in A$ 表示 $x$ 是集合 $A$ 的元素；$A\subseteq B$ 表示 $A$ 的每个元素都属于 $B$。若 $A\subseteq B$ 且 $B\subseteq A$，则称两集合相等。没有元素的集合记为 $\varnothing$。
\end{definition}
集合不记录排列次序和重复次数，故 $\{1,2,2\}=\{2,1\}$。描述法写成 $A=\{x\in D:P(x)\}$，其中 $D$ 是对象范围，$P(x)$ 是筛选条件。证明两个集合相等，通常分别证明两方向的包含。空集是每个集合的子集：若这一结论失败，就必须在空集中找到一个不属于目标集合的元素，而空集没有元素。

实数区间 $(a,b)=\{x:a<x<b\}$ 不含端点，$[a,b]=\{x:a\le x\le b\}$ 包含端点；半开区间相应处理。符号 $\pm\infty$ 不表示实数端点，因此区间在无穷一侧总用圆括号。必须区分元素关系 $1\in\{1,2\}$ 与包含关系 $\{1\}\subseteq\{1,2\}$。

\section{集合运算与分类表达}
在固定全集 $U$ 中，定义并集 $A\cup B=\{x:x\in A\text{ 或 }x\in B\}$、交集 $A\cap B=\{x:x\in A\text{ 且 }x\in B\}$、差集 $A\setminus B=\{x:x\in A,\ x\notin B\}$ 和补集 $A^c=U\setminus A$。这里“或”允许两者同时成立；“恰有一个成立”对应 $(A\setminus B)\cup(B\setminus A)$。

\begin{proposition}[德摩根律]
对 $U$ 的任意子集 $A,B$，有 $(A\cup B)^c=A^c\cap B^c$ 及 $(A\cap B)^c=A^c\cup B^c$。
\end{proposition}
\begin{proof}
元素 $x$ 不属于 $A\cup B$，当且仅当它既不属于 $A$ 也不属于 $B$，这正是属于 $A^c\cap B^c$ 的条件。第二式中“不同时属于两者”等价于“至少不属于其中之一”，同理得到结论。
\end{proof}
把一个集合写成若干子集的并称为分类；若各部分两两不交且覆盖原集合，称为分割。分类讨论必须覆盖全部情形；计数时若各类有交叠，还需修正重复。

\section{命题、量词与条件关系}
\emph{命题}（proposition）是具有确定真假值的陈述。$P\Rightarrow Q$ 表示只要 $P$ 成立，$Q$ 就成立，称 $P$ 是 $Q$ 的充分条件、$Q$ 是 $P$ 的必要条件；双向均成立写为 $P\Leftrightarrow Q$。例如对实数 $x$，$x>0$ 能推出 $x^2>0$，逆向不成立；正确的等价关系是 $x^2>0\Leftrightarrow x\ne0$。逆否命题 $\neg Q\Rightarrow\neg P$ 与原命题等价，逆命题 $Q\Rightarrow P$ 一般不等价。

全称量词 $\forall$ 表示“对所有”，存在量词 $\exists$ 表示“至少有一个”。否定时既要否定谓词，也要交换量词：
\[
 \neg(\forall x\in D,\ P(x))\Longleftrightarrow
 \exists x\in D,\ \neg P(x),\qquad
 \neg(\exists x\in D,\ P(x))\Longleftrightarrow
 \forall x\in D,\ \neg P(x).
\]
“对每个实数 $x$，存在实数 $y>x$”为真，可取 $y=x+1$；交换量词成为“存在实数 $y$ 大于每个实数 $x$”则为假，取 $x=y$ 即违反要求。变量的依赖关系因此是命题的一部分。

\section{直接证明、反证与反例}
直接证明从假设与定义出发。例如偶数 $a=2m,b=2n$ 的和为 $2(m+n)$，仍为偶数；这里 $m,n\in\mathbb Z$ 保证 $m+n$ 是整数。证明“整数 $n$ 的平方为偶数则 $n$ 为偶数”，可证明逆否命题：若 $n=2k+1$，则 $n^2=2(2k^2+2k)+1$ 为奇数。

\emph{反证法}（proof by contradiction）假定结论不成立，并从这一假定与原条件推出矛盾。它不能偷偷改动原条件，也不能把待证结论作为推理依据。否定全称命题只需一个满足全部假设却违反结论的反例，例如 $x=-1$ 否定了“$x^2=1$ 必有 $x=1$”；验证许多个正例仍不能证明该命题对所有对象成立。

\section{数学归纳法与递归描述}
\begin{theorem}[数学归纳法]
设 $P(n)$ 是对整数 $n\ge n_0$ 定义的命题。若 $P(n_0)$ 成立，且对每个 $n\ge n_0$ 都有 $P(n)\Rightarrow P(n+1)$，则所有 $n\ge n_0$ 上的 $P(n)$ 均成立。
\end{theorem}
\begin{proof}
采用正整数的良序性质，即非空正整数集合有最小元素。若存在反例，取最小反例 $m$。起点已验证，故 $m>n_0$；于是 $P(m-1)$ 成立，归纳步骤又推出 $P(m)$，矛盾。对任意整数起点，可先平移指标到正整数。
\end{proof}
例如 $1+\cdots+n=n(n+1)/2$：$n=1$ 时成立；若对 $n$ 成立，则增加 $n+1$ 后得到 $(n+1)(n+2)/2$，完成归纳。归纳假设并未假定所有情形成立，只把某一步成立作为推出下一步的条件。强归纳法允许在推导 $P(n+1)$ 时使用从起点到 $n$ 的全部结论，证明同样依赖最小反例。

递归定义则规定对象如何逐步生成，如 $a_1=1,\ a_{n+1}=a_n+2$ 唯一确定正奇数列。初值不可省略，否则许多不同数列可能遵守同一更新规则。

\section{例题、习题与学习检查}
\begin{example}
设 $A=\{x\in\mathbb R:x^2=1\}$，$B=\{x\in\mathbb R:x>0\}$。求同时或恰有一个条件成立的实数。
\end{example}
\begin{solve}
$A=\{-1,1\}$，故同时成立的集合为 $A\cap B=\{1\}$；恰有一个成立的集合为 $\{-1\}\cup((0,+\infty)\setminus\{1\})$。这里 $1$ 虽属于并集，却应从“恰有一个”的答案中剔除。
\end{solve}
\begin{enumerate}
\item 否定命题“对每个正整数 $n$，存在正整数 $m>n$ 且 $m$ 为偶数”。答：存在正整数 $n$，使每个正整数 $m$ 都满足 $m\le n$ 或 $m$ 为奇数。
\item 证明 $1+3+\cdots+(2n-1)=n^2$。提示：归纳步骤增加 $2n+1$。
\item 判断 $A\subseteq B\cup C$ 是否必有 $A\subseteq B$ 或 $A\subseteq C$。反例为 $A=\{1,2\},B=\{1\},C=\{2\}$。
\end{enumerate}

